\chapter{因果识别与管理干预}
\label{ch:week08-causality}

因果分析研究改变处理状态对结果的影响。加急订单的延期率可能高于普通订单，
因为管理者优先处理困难订单。评价加急效果，需要比较同类订单接受加急与保持
原安排时的结果。实际只观察其中一种状态，另一种状态的结果构成反事实。

因果分析须区分目标效应、识别条件与样本估计量。标准化与加权依赖给定协变量
的可比性，差分中的差分依赖未处理结果的平行趋势。双重稳健估计结合结果回归
与处理概率模型。观察到的组间差异为何偏离目标效应，需要结合处理的选择机制分析。

\section{干预效果的研究问题}

发券对象通常经过选择，其风险、价值、活跃度与人工服务情况可能不同于未发券对象。两组留存率的差异因此同时反映处理效果和原有人群差异，需要进一步分析分配机制。

\begin{definition}[因果问题]
因果问题关注行动或处理 \(D\) 对结果 \(Y\) 的影响，以同一目标人群在不同处理状态下的潜在结果差异定义效应。观察组间均值差能否识别该效应，取决于处理分配机制与相关假设。
\end{definition}

管理场景中的因果问题包括：优惠券是否提高留存，新排班规则是否减少延迟，价格调整是否提升毛利，推荐策略是否增加复购，供应商整改是否改善准时率，审批自动化是否缩短周期。它们的共同难点是，组织只能观察对象实际接受的行动及其结果，不能同时观察同一对象在另一种行动下的结果。

\section{反事实、偏差与识别}
\subjectindex{yinguoshibie}{因果识别}
\subjectindex{fanshishi}{反事实}
\subsection{潜在结果与反事实}
\subjectindex{ATE}{ATE（平均处理效应）}
\subjectindex{ATT}{ATT（已处理者平均效应）}

对个体 \(i\)，令 \(D_i\in\{0,1\}\) 表示是否接受处理。潜在结果记为
\[
  Y_i(1):\text{接受处理时的结果},
  \qquad
  Y_i(0):\text{不接受处理时的结果}.
\]
个体处理效应为
\[
  \tau_i=Y_i(1)-Y_i(0).
\]
但是观察到的结果只有
\[
  Y_i=D_iY_i(1)+(1-D_i)Y_i(0).
\]
若客户收到了优惠券，只能观察发券后的结果 \(Y_i(1)\)，没有发券时的结果 \(Y_i(0)\) 是反事实；若客户未收到优惠券，则发券后的结果是反事实。这就是 Holland 所称的基本因果问题：同一个体的两个潜在结果不能被同时观察。

\begin{definition}[平均处理效应]
总体平均处理效应为
\[
  ATE=\E[Y(1)-Y(0)].
\]
若只关心实际接受处理的人群，则处理组平均处理效应为
\[
  ATT=\E[Y(1)-Y(0)\mid D=1].
\]
若关心特定协变量取值 \(X=x\) 的人群，则条件平均处理效应为
\[
  CATE(x)=\E[Y(1)-Y(0)\mid X=x].
\]
\end{definition}

估计对象应由管理问题确定。评价在目标人群中全面实施优惠券策略的平均效果，可关注 ATE；评价实际被运营团队触达客户的平均效果，可关注 ATT；选择行动对象时，则需结合 CATE、收益与行动成本建立优化问题，见第~\ref{ch:week09-optimization}~章。

\subsection{统计关联、预测与因果效应}

相关分析、预测、因果推断和处方分析具有不同的目标及解释边界。

\begin{table}[htbp]
\centering
\caption{几类分析问题的区别}
\begin{tabular}{P{0.18\linewidth}P{0.36\linewidth}P{0.34\linewidth}}
\toprule
\rowcolor{BDMTableHead}
任务 & 典型问题 & 主要风险\\
\midrule

相关分析 & 客服工单数是否与流失率同向变化？ & 把共同变化解释为行动效果。\\
预测建模 & 哪些客户未来最可能流失？ & 把预测特征解释为可干预原因。\\
因果推断 & 发放优惠券是否提高留存？ & 忽视混杂、选择偏差和识别假设。\\
处方分析 & 在预算约束下应对哪些客户行动？ & 用未经验证的效果估计优化资源。\\
\bottomrule
\end{tabular}
\end{table}

例如，折扣率与流失风险可能正相关，因为高风险客户更常获得折扣；这不意味着提高折扣会导致流失。预测模型可以利用折扣率提升排序能力，但因果报告不能把该系数解释为行动效果。报告应按研究问题区分描述结果、预测结果、因果效应及决策建议。

\subsection{选择偏差与混杂}
\subjectindex{hunza}{混杂}

朴素均值差可以写成
\[
  \E[Y\mid D=1]-\E[Y\mid D=0].
\]
代入潜在结果并整理，可得到
\[
\begin{aligned}
  \E[Y\mid D=1]-\E[Y\mid D=0]
  &=
  \E[Y(1)-Y(0)\mid D=1] \\
  &\quad +
  \{\E[Y(0)\mid D=1]-\E[Y(0)\mid D=0]\}.
\end{aligned}
\]
第一项是 ATT，第二项是选择偏差。若处理组在没有处理时本来也会有不同结果，朴素差异就不等于处理效应。

\begin{definition}[混杂变量]
混杂变量是同时影响处理分配 \(D\) 和结果 \(Y\) 的变量。若风险分数既影响运营人员是否发券，也影响客户留存，那么风险分数就是优惠券效果评估中的混杂变量。
\end{definition}

混杂不同于中介变量和碰撞变量。中介变量位于处理之后，是处理影响结果的路径之一；例如优惠券改变客户使用活跃度，再影响留存。若估计总效果，通常不应控制处理后的中介。碰撞变量则是由两个变量共同影响的变量；错误控制碰撞变量可能引入原本不存在的关联。管理数据中常见的“进入人工复核队列”就可能是碰撞变量，因为它同时受风险评分和客户经理判断影响。

\subsection{混杂变量、中介变量与碰撞变量}
\subjectindex{yinguotu}{因果图}
\subjectindex{zhongjie}{中介}
\subjectindex{pengzhuangdian}{碰撞点}
因果图的箭头表示研究假定的直接因果关系。若设备老化 \(X\) 同时影响是否维修 \(D\) 与故障 \(Y\)，路径 \(D\leftarrow X\rightarrow Y\) 是需要处理的混杂路径。若维修先改变温度 \(M\)，温度再影响故障，\(M\) 是中介；控制它会改变总效应问题。若维修与潜在故障都会触发复核 \(C\)，只分析被复核样本可能打开碰撞路径。

\begin{figure}[htbp]
\centering
\begin{tikzpicture}[every node/.style={font=\small}]
\node[bdmvar] (x) at (0,1.4) {\(X\)};
\node[bdmvar] (d) at (-1.3,0) {\(D\)};
\node[bdmvar] (y) at (1.3,0) {\(Y\)};
\draw[bdmarrow] (x)--(d);\draw[bdmarrow] (x)--(y);\draw[bdmarrow] (d)--(y);
\node[align=center,text=BDMGray] at (0,-1) {（a）混杂\\调整前置共同原因};
\node[bdmvar] (d2) at (3.3,0) {\(D\)};
\node[bdmvar] (m2) at (4.6,0) {\(M\)};
\node[bdmvar] (y2) at (5.9,0) {\(Y\)};
\draw[bdmarrow] (d2)--(m2);\draw[bdmarrow] (m2)--(y2);
\node[align=center,text=BDMGray] at (4.6,-1) {（b）中介\\控制后不再是总效应};
\node[bdmvar] (d3) at (7.9,0) {\(D\)};
\node[bdmvar] (c3) at (9.2,1.4) {\(C\)};
\node[bdmvar] (y3) at (10.5,0) {\(Y\)};
\draw[bdmarrow] (d3)--(c3);\draw[bdmarrow] (y3)--(c3);
\node[align=center,text=BDMGray] at (9.2,-1) {（c）碰撞\\条件化可能引入关联};
\end{tikzpicture}
\caption{三个最小因果结构。图是业务机制假设，需要知识与证据支持。}
\end{figure}

智能体可以根据业务描述提出候选因果图，但仅凭观察到的相关性，通常不能唯一确定真实因果图。主管未记录的紧急程度、设备健康和客户价值尤其可能造成遗漏混杂。应把每条箭头的依据与争议写入研究设计。

\subsection{随机实验作为识别基准}

若处理 \(D_i\) 由随机机制决定，并且实验执行稳定，则
\[
  D_i\perp (Y_i(1),Y_i(0)).
\]
处理组与对照组在期望上可比，因此
\[
  \E[Y\mid D=1]-\E[Y\mid D=0]
  =
  \E[Y(1)-Y(0)].
\]
这正是第~\ref{ch:week06-ab-testing}~章 A/B 测试的理论基础。随机化并不保证每次小样本都完全平衡，但它在设计上切断了处理分配与潜在结果之间的系统性联系。

在企业环境中，随机实验还要求处理定义清晰、执行稳定、样本比例正常、结果窗口一致，并尽量避免实验污染。例如客户之间共享优惠券、销售团队绕开分组规则、推荐系统同时改变多个模块，都会削弱随机化带来的识别优势。

\begin{modernbox}[设计优先]
因果识别依赖处理分配机制、明确的干预定义和可观察的结果。估计模型不能弥补关键识别条件的缺失。
\end{modernbox}

\subsection{观测性研究的识别假设}
\subjectindex{chongdiexing}{重叠性}
\subjectindex{wuhunza}{无混杂}
\label{sec:causal-assumptions}

许多管理干预不能完全随机化。例如，服务承诺可能不允许对高价值客户随机拒绝服务，安全要求可能不允许对高风险订单随机停止处置，价格政策也可能按区域逐步上线。此时需要在观测性数据中构造可比性。常见识别假设如下。

\begin{definition}[一致性]
若个体实际接受处理 \(D_i=d\)，则观察到的结果等于对应潜在结果：
\[
  Y_i=Y_i(d).
\]
这要求处理定义足够清晰。例如“发优惠券”必须说明金额、渠道、有效期、触达时间和是否允许叠加。
\end{definition}

\begin{definition}[SUTVA]
SUTVA 要求一个个体的潜在结果不受其他个体处理状态影响，且处理没有隐藏版本。若客户之间会转赠优惠券，或不同客户经理执行方式差异很大，SUTVA 可能失败。
\end{definition}

\begin{definition}[条件可交换性]
给定可观察协变量 \(X\) 后，处理分配与潜在结果独立：
\[
  (Y(1),Y(0))\perp D\mid X.
\]
这意味着所有同时影响处理和结果的重要混杂因素都已被观察并纳入调整。
\end{definition}

\begin{definition}[重叠性]
对每个协变量取值 \(x\)，处理概率不能为 0 或 1：
\[
  0<\Prob(D=1\mid X=x)<1.
\]
若某类客户总是被发券，而另一类客户从不发券，就无法仅凭该协变量层内的观测比较两种处理状态。
\end{definition}

在一致性、条件可交换性和重叠性下，条件结果均值识别潜在结果均值。
对 \(d\in\{0,1\}\)，有
\[
\E[Y(d)\mid X=x]
=\E[Y(d)\mid D=d,X=x]
=\E[Y\mid D=d,X=x].
\]
第一步使用条件可交换性，第二步使用一致性；重叠性保证两种处理下的条件均值均有
观测依据。再对目标总体的 \(X\) 分布取平均，即得标准化识别公式
\[
\operatorname{ATE}
=\E_X\{m_1(X)-m_0(X)\},\qquad
m_d(x)=\E[Y\mid D=d,X=x].
\]

\subsection{混杂调整与标准化算例}
\label{sec:causal-standardization}
\subjectindex{biaozhunhua_yinguo}{标准化（因果推断）}
考虑100张订单，其中50张为高难度、50张为低难度。加急标为 \(D=1\)，延期标为
\(Y=1\)。表~\ref{tab:causal-example}~给出一组教学数据。
\begin{table}[htbp]
\centering
\caption{困难对象优先处理时的选择偏差}
\label{tab:causal-example}
\begin{tabular}{lrrrr}
\toprule
\rowcolor{BDMTableHead}
难度 & 加急订单数 & 加急延期数 & 未加急订单数 & 未加急延期数\\
\midrule
高 & 40 & 12 & 10 & 4\\
低 & 10 & 0 & 40 & 4\\
\bottomrule
\end{tabular}
\end{table}

直接比较得到加急组延期率 \(12/50=0.24\)，未加急组为 \(8/50=0.16\)，
差异为0.08。这个差异描述两组已观察结果，却没有控制加急组包含更多困难订单的事实。

在相同难度内比较，高难度层的差异为 \(12/40-4/10=-0.10\)，低难度层
为 \(0/10-4/40=-0.10\)。若目标总体高、低难度各占一半，
将两个行动分别标准化到同一构成，得到
\[
\widehat\mu_1=0.5(0.30)+0.5(0)=0.15,\qquad
\widehat\mu_0=0.5(0.40)+0.5(0.10)=0.25.
\]
标准化差异为 \(-0.10\)，即在固定难度构成下，加急组延期率低10个百分点。

标准化结果的因果解释还需要识别条件。如果难度已经包括所有同时影响加急
选择与延期的处理前因素，并且每一层都有两种行动、结果定义一致、不存在相关干扰，
上述标准化才具有相应因果解释。若主管还依据未记录的设备故障选择加急，分层之后
仍可能混杂，因而需要结合处理分配机制论证调整变量是否充分。

标准化与第~\ref{ch:week05-statistics-dashboard}~章的组合分解具有相近的加权结构。
组合分解描述权重变化造成的总体差异；因果标准化还要求层内对象满足反事实比较
的识别条件。

\section{倾向得分与加权估计}
\subsection{分层、匹配与倾向得分}
\subjectindex{qingxiangdefen}{倾向得分}

若混杂变量 \(X\) 可观察，一个直接办法是在相似 \(X\) 的对象之间比较处理和对照。分层法把样本划分为若干层，在层内比较均值差，再按层权重合成总体估计。分层结果易于核查，但协变量维度较高时，部分层内的样本可能很少，甚至缺少某一处理组。

倾向得分是给定协变量后的处理概率：
\[
  e(X)=\Prob(D=1\mid X).
\]
Rosenbaum 与 Rubin 的结果包含平衡性质与因果识别两个层次。真实倾向得分首先是平衡得分，即
\(D\perp X\mid e(X)\)；这一性质来自倾向得分的定义，不要求无混杂。
若进一步满足给定 \(X\) 的条件可交换性与重叠性，才能将潜在结果的可交换性
转移到给定 \(e(X)\) 的比较中。估计得分相近并不保证有限样本中的协变量已经平衡，也不能
证明未观测因素可比；仍须进行平衡诊断并分析业务机制。

倾向得分的常见使用方式包括：
\begin{enumerate}
  \item 分层：按倾向得分或关键风险分数分层，在层内比较处理和对照；
  \item 匹配：为每个处理对象寻找倾向得分相近的对照对象；
  \item 逆概率加权：处理组权重为 \(1/e(X)\)，对照组权重为 \(1/(1-e(X))\)；
  \item 平衡诊断：比较调整前后协变量分布是否更接近。
\end{enumerate}

倾向得分模型估计处理概率，选择模型时应检查概率估计、协变量平衡、重叠范围和稳定性。接近0或1的估计可能反映真实的处理规则，也可能来自过拟合。截断权重或缩小目标人群时，应说明相应的偏差和估计对象变化。

\subsection{逆概率加权与平衡诊断}
\subjectindex{IPW}{IPW（逆概率加权）}
\subjectindex{pinghengzhenduan}{平衡诊断}

逆概率加权估计 ATE 的一个形式为
\[
  \widehat{ATE}_{IPW}
  =
  \frac{1}{n}\sum_{i=1}^{n}
  \left[
    \frac{D_iY_i}{\hat e(X_i)}
    -
    \frac{(1-D_i)Y_i}{1-\hat e(X_i)}
  \right].
\]
其依据是条件期望恒等式。在真实倾向得分下，
\[
\E\!\left[\frac{DY}{e(X)}\mid X\right]
=\frac{e(X)\E[Y\mid D=1,X]}{e(X)}=m_1(X).
\]
对照项同理。对 \(X\) 取期望后相减，便得到前述标准化公式。样本估计以
\(\hat e(X)\) 代替未知概率；较少接受某处理的对象在实际接受该处理时获得较大权重。

平衡诊断常用标准化均值差异：
\[
  SMD(X_j)
  =
  \frac{\bar X_{j,T}-\bar X_{j,C}}
       {\sqrt{(s_{j,T}^{2}+s_{j,C}^{2})/2}}.
\]
加权后重新计算 SMD，可检查观测协变量的分布差异。残余不平衡提示需要调整模型或限制比较范围；未观测混杂则需结合领域知识、研究设计和敏感性分析\footnote{因果分析中的敏感性分析考察识别假设偏离到何种程度会改变结论，例如设定遗漏混杂可能造成的偏差范围。它不能证明没有遗漏混杂，而是说明结论对该风险的敏感程度。}考察。

\subsection{有效样本量与重叠性}
\subjectindex{ESS}{ESS（有效样本量）}
对一组非负权重 \(w_i\)，常用诊断量为
\[
n_{\rm eff}=\frac{(\sum_iw_i)^2}{\sum_iw_i^2}.
\]
该式要求至少一个权重为正；全部权重为零时不可定义。全部权重相等且为正时，
它等于观测数；少数权重很大时明显下降。该量描述权重集中程度，估计量的标准误还取决于结果变量及抽样结构。应分处理组和对照组检查倾向得分分布、权重最大值及加权后协变量平衡。

截断极端权重可能降低方差，同时引入偏差或改变目标人群。删除没有重叠的高风险批次，意味着结果只适用于剩余可比范围。报告应据此重新定义目标总体及其平均处理效应。

\section{双重稳健估计与因果机器学习}

\subsection{增广逆概率加权估计}
\subjectindex{AIPW}{AIPW}
\subjectindex{shuangchongwenjian}{双重稳健}
定义 \(m_d(x)=\E[Y\mid D=d,X=x]\) 和 \(e(x)=\Prob(D=1\mid X=x)\)。在一致性、条件可交换性、重叠性及适当独立抽样条件下，ATE 的增广逆概率加权估计为
\[
\hat\tau=\frac1n\sum_i\left[
\hat m_1(X_i)-\hat m_0(X_i)
+\frac{D_i(Y_i-\hat m_1(X_i))}{\hat e(X_i)}
-\frac{(1-D_i)(Y_i-\hat m_0(X_i))}{1-\hat e(X_i)}
\right].
\]
第一部分使用结果模型预测两种状态，后两部分分别用处理组和对照组的加权残差修正预测。该估计量同时利用结果回归和倾向得分。

\subsection{双重稳健性的条件期望证明}
先假设倾向得分正确，而结果模型可能错误。将两个候选结果函数写为 \(\tilde m_d\)。给定 \(X=x\)，
\[
\E\!\left[\frac{D(Y-\tilde m_1(x))}{e(x)}\mid X=x\right]
=m_1(x)-\tilde m_1(x).
\]
对照项同理，两项校正恰好抵消结果模型的偏差。反过来，若结果模型正确，无论候选倾向得分取何种严格位于 0 与 1 之间的值，残差项条件期望均为零，剩下 \(m_1-m_0\)。两种情形均可消除总体分数的偏差。

在前述识别假设与适当正则条件下，倾向得分模型正确，或两个结果回归模型均正确，
可分别支持该估计量的一致性。这一性质称为双重稳健性。

也可以把偏差明确写出。令 \(\phi\) 表示采用固定候选函数 \(\tilde e,\tilde m_0,\tilde m_1\) 的单样本分数，在给定 \(X=x\) 后省略函数参数，有
\[
\E[\phi\mid X=x]-(m_1-m_0)
=(\tilde e-e)\left[
\frac{\tilde m_1-m_1}{\tilde e}
+\frac{\tilde m_0-m_0}{1-\tilde e}\right].
\]
因此，倾向得分误差为零，或两种结果回归误差均为零，都会使条件偏差消失；当两类误差都较小时，它们以乘积而非简单相加进入偏差。候选倾向得分接近0或1时，分母会放大两类估计误差的影响。

\subsection{正交性与交叉拟合}
\subjectindex{jiaochanihe}{交叉拟合}
若用同一观测拟合复杂模型并计算其残差，过拟合会破坏简单的渐近推断。交叉拟合先把数据分成 \(K\) 折，在其余折估计 \(\hat e,\hat m_d\)，再只在留出折计算上述分数，最后汇总。因此，计算每个样本的分数时，所用辅助模型均未使用该样本拟合。

双重机器学习\footnote{双重机器学习（double/debiased machine learning，DML）利用机器学习估计结果回归、倾向得分等辅助函数，再结合正交估计方程与样本分割推断目标效应；“双重”不表示任意两个模型拟合良好就能保证因果识别。}利用估计方程对辅助函数扰动的一阶不敏感性，即局部正交性，使小的
估计误差主要通过乘积影响偏差\cite{dml08}。记
\(\|h\|_2=(\E[h(X)^2])^{1/2}\)。在倾向得分有界远离0和1、适当矩条件及辅助函数
一致收敛等条件下，一个常用的速率要求是
\[
\|\hat e-e\|_2\bigl(\|\hat m_1-m_1\|_2+\|\hat m_0-m_0\|_2\bigr)
=o_p(n^{-1/2}).
\]
记号 \(Z_n=o_p(a_n)\) 表示 \(Z_n/a_n\) 依概率趋于0。若两类误差均为
\(o_p(n^{-1/4})\)，则满足该乘积条件。交叉拟合负责分离拟合与评价样本，收敛速率
则取决于数据分布与所用估计方法。

\begin{example}[单个双重稳健分数]
某样本 \(D=1,Y=1,\hat e=0.5,\hat m_1=0.8,\hat m_0=0.6\)，其分数为 \(0.8-0.6+(1-0.8)/0.5=0.6\)。另一个对照样本 \(D=0,Y=0\)，使用同样预测得到 \(0.2-(0-0.6)/0.5=1.4\)。单个分数可以超过二元结果差的范围；该分数是用于总体平均估计的加权校正量。
\end{example}

\section{差分中的差分与政策评价}
\subsection{差分中的差分}
\subjectindex{DID}{DID（差分中的差分）}
\subjectindex{pingxingqushi}{平行趋势}

有些干预按时间逐步实施。例如部分门店上线新的留存策略，其他门店暂未上线。若能观察处理组和对照组在干预前后的结果，可以使用差分中的差分，简称 DID。

设 \(Y_{gt}\) 表示组 \(g\) 在时期 \(t\) 的平均结果。两期 DID 估计为
\[
  \widehat{\tau}_{DID}
  =
  (\bar Y_{T,post}-\bar Y_{T,pre})
  -
  (\bar Y_{C,post}-\bar Y_{C,pre}).
\]
第一项是处理组的前后变化，第二项是对照组的前后变化。DID 的思想是用对照组的同期变化近似处理组在没有干预时的趋势。

\begin{definition}[平行趋势假设]
若没有干预，处理组和对照组的平均结果变化趋势应相同。形式上，
\[
  \E[Y_{T,post}(0)-Y_{T,pre}(0)]
  =
  \E[Y_{C,post}(0)-Y_{C,pre}(0)].
\]
\end{definition}

两期 DID 还要求处理对干预前结果没有预期影响，对照组在研究期内未接受相应处理，且两组的比较对象保持一致。在这些条件下，将处理组事后结果分解为未处理结果与处理效应，再利用平行趋势抵消两组的未处理变化，即可识别处理组的平均效应。干预前多期趋势、业务机制与安慰剂检验\footnote{安慰剂检验将同一方法用于理论上不应受处理影响的时期、对象或结果，检查是否仍得到明显偏离零的效应估计。检验通过可增加设计的可信度，但不能证明干预后的平行趋势成立。}可提供经验依据；同期政策和组别构成变化则需另外分析。

\subsection{分期处理与异质性效应}
\subjectindex{fenqichuli}{分期处理}
假设产线甲在一月采用新维护策略，乙在三月采用，丙尚未采用。若效应随采用时长或产线不同，简单双向固定效应回归\footnote{双向固定效应回归同时控制组别固定差异与各组共同的时期冲击。分期处理且效应异质时，其回归系数未必是目标处理效应的适当加权平均。}可能把已经处理的甲当作乙的对照，混入甲已有的处理效应。Sun 与 Abraham、Callaway 与 Sant'Anna 的研究明确讨论了分期处理与异质效应问题\cite{sun08,csdid08}。

可定义首次处理时点 \(G=g\) 的组在时期 \(t\) 的
\[
\operatorname{ATT}(g,t)=\E[Y_t(g)-Y_t(\infty)\mid G=g],\quad t\ge g,
\]
其中 \(Y_t(g)\) 是首次在 \(g\) 时点接受处理后于 \(t\) 的结果，\(Y_t(\infty)\) 是始终未处理的结果；该定义保留首次处理时间及处理持续时间。使用尚未处理或从未处理、且满足相应平行趋势条件的对象构造对照，再按明示的组别权重汇总。还需考虑预期效应、群体构成变化和同期政策。干预前趋势提供部分经验依据，干预后的平行趋势仍需结合业务机制论证。

\section{处理效应与决策规则}

\subsection{处理时点与研究设计}
研究设计应明确资格判定、可选行动、处理开始和随访结果的时点，使处理分组与随访起点一致。若使用“未来是否完成
加急流程”来划分处理组，却从更早日期开始计算结果，就可能让某一组被要求先存活
或保持未结案一段时间才能进入分析。这种按未来状态分组引起的不死时间偏倚，可通过对齐资格判定、处理开始和随访起点加以避免。

协变量的选择取决于其在因果结构中的作用。预测能力强的变量未必适合调整；
行动之后才出现的设备温度、客户投诉或返工记录可能是中介，也可能包含碰撞点。
协变量集合应由因果图及时间顺序确定，机器学习随后用于估计这些变量之间的条件关系。

敏感性分析可用偏差幅度考察不可检验的识别假设。例如，若假定估计结果相对于
目标效应仍可能有绝对值不超过 \(b\) 的未测混杂偏差，那么一个
原统计区间 \([l,u]\) 可保守扩展为 \([l-b,u+b]\)，前提是原区间覆盖其统计目标、
偏差界本身成立。\(b\) 是敏感性分析预先设定的偏差幅度，可用于考察多大的遗漏影响会改变结论。若很小的 \(b\) 就能推翻建议，应降低行动强度或设计更可信的试点。

\subsection{效应异质性与外部有效性}

相同的平均处理效应可能对应不同的人群效应分布。优惠券可能对高风险客户有效，对低风险客户无效；对高价值客户值得投放，对低价值客户即使有效也不划算。行动选择需要结合效应异质性、实施时点与成本。

因果估计的外推必须谨慎。若样本只包含某一区域、某一渠道或某一风险层，估计结果未必适用于所有客户。若重叠性只在中等风险客户中成立，模型就不能可靠回答极高风险客户的反事实。报告应明确估计适用的人群及向其他人群推广所需的条件。

\subsection{策略价值与预算}
\subjectindex{celuejiazhi}{策略价值}
若 \(Y=1\) 表示保留客户，定义条件处理效应 \(\tau(x)=\E[Y(1)-Y(0)\mid X=x]\)，并令 \(\pi(x)\in\{0,1\}\) 为是否发券的策略。相对全部不处理的净价值为
\[
V(\pi)=\E[\pi(X)\{v(X)\tau(X)-c(X)\}],
\]
约束可以是 \(\E[\pi(X)c(X)]\le B\)。当成本相同且只有人数上限时，可按正的净增益排序；成本不同且需整数选择时，应进入第~\ref{ch:week09-optimization}~章的背包或混合整数模型。按净收益成本比排序并不总能求得整数问题最优解。

策略学习研究如何从观测数据学习此类规则，并控制相对最优策略的福利损失\cite{policy08}。训练效应模型、选择策略和评估策略不能全部使用同一数据而不作校正。高度灵活的策略可能因样本噪声而选中效应被高估的个体，使同一样本上的收益评价偏高。

\section{数据平台与人工智能的效果评价}

\subsection{企业平台应用的识别设计}
“采用数据平台后交付率提高”包含至少三种解释：平台有贡献，采用部门本来就在改善，或同期更换了设备与管理制度。Palantir、Oracle、Atlan、Collibra 等公开案例可以说明应用路径，但若要估计净经营效应，应进一步取得采用时点、可比对象、结果口径和共同冲击信息。

考虑“某制造集团分批部署统一数据目录”的假设案例。采用时点可定义为业务团队开始使用指定指标服务的日期。结果可以是完成标准分析任务的总工时、复核错误率与争议解决时间；登录量只能反映使用情况，不能单独代表上述经营效果。采用顺序若优先选择准备程度高的团队，便存在选择偏差。可随机安排培训名额、在可行条件下随机部署次序，或明确采用观测性设计及其假设。

\subsection{AI 使用量、个人节时与组织生产率}
平台日志通常只覆盖使用该平台的人。Anthropic 2026 年 Economic Index 的日志和调查材料适合研究采用与任务模式，但不能单独识别没有使用 AI 时的反事实生产率\cite{aiindex08}。高技能员工可能更早采用，困难任务也可能更常求助于 AI，这两种选择机制对观察到的生产率差异可能产生相反方向的偏差。

AI 应用评估应记录任务难度、首次完成时间、人工复核、返工、质量和最终业务结果。若采用 AI 的团队减少了编写时间却增加了其他部门审核时间，个人效率提高可能没有转成组织收益。研究应先固定评价单位为任务、员工、团队还是业务流程，再选择相应的随机化或比较设计。

\subsection{平台组合的经营效果}
一次数智化改造可能同时更换数据目录、查询系统、模型端点与工单流程。
此时“使用 AI 平台”的处理包含多个共同变化的组成部分。若目标是评估整体部署，
应把完整配置及上线时点定义为处理；若目标是比较语义服务的增量价值，应在两组
保持其他组件一致的条件下改变语义服务。两种问题对应不同的反事实，
仅凭同一个上线前后差值无法分别识别。

可将部门 \(g\) 在时段 \(t\) 的每百张到期订单误判数记为 \(Y_{gt}\)。按每百张到期订单
计算误判数，可使处理组与对照组在订单量不同的时期具有可比的统计尺度；还要检查
客户结构和交期定义是否变化。若平台按部门分期上线，可依据第4节的差分中的
差分方法定义各采用批次的处理效应，核查上线前趋势、同期制度变化以及跨部门
共享人员造成的干扰。对统一服务中心而言，员工在多部门轮转可能使尚未上线
部门间接受到处理，这一事实应进入设计记录。

经济评价把减少的误判损失、节省的净工时与新增的模型、工具、复核和平台费用
置于同一观察期。若以 \(V\) 表示在因果研究条件下估计的经营收益，\(F\) 为
一次性实施成本，\(C_t\) 为第 \(t\) 期运行成本，贴现率 \(r\ge0\)，则
\[
  \operatorname{NPV}=V-F-\sum_{t=1}^{T}\frac{C_t}{(1+r)^t},
\]
其中 \(V\) 已按同一估计期和贴现口径计算。若只有效果的上下界，应相应报告
净现值区间；引用供应商案例只能说明产品机制，不能替代 \(V\) 的识别。

\subsection{因果研究的实施步骤}
目标试验是研究者希望用来回答因果问题的理想随机试验。分析观测数据时，先按这一试验规定资格条件、处理方案、分配时点、随访期和结果变量，再考察现有资料能够对应哪些要素\cite{hernan2020}。再画出处理分配机制，区分前置变量与处理后变量。随后检查重叠、缺失和时间，选择与识别策略一致的估计方法。最后报告效应区间、敏感性和外推范围。人工智能工具可辅助检查报告要素，未观测混杂的分析依据业务机制和敏感性设定。

\section{观测性数据的计算实例}

本节计算分层、逆概率加权、平衡诊断与两期差分中的差分。两个算例分别说明横截面调整和时间比较，使用各自明确的处理与结果定义。

本章有两份样例数据。第一份为观测性客户优惠券数据：
\begin{center}
\path{data/sample/coupon_retention_observational.csv}
\end{center}
第二份为门店层面的两期面板数据：
\begin{center}
\path{data/sample/coupon_retention_panel.csv}
\end{center}
配套代码位于：
\begin{center}
\path{src/bdm_decision/cases/week08_causality.py}
\end{center}

观测性优惠券数据包含 24 个客户，其中 12 个收到优惠券，12 个没有收到。由于运营人员更倾向于给高风险客户发券，处理组和对照组并不直接可比。脚本得到：
\begin{table}[htbp]
\centering
\caption{教学优惠券数据的三种留存效应估计}
\begin{tabular}{lrrr}
\toprule
\rowcolor{BDMTableHead}
方法 & 处理组均值 & 对照组均值 & 效果估计\\
\midrule
朴素差异 & 0.7500 & 0.7500 & 0.0000\\
风险分层 & NA & NA & 0.2500\\
逆概率加权 & NA & NA & 0.2199\\
\bottomrule
\end{tabular}
\end{table}

分层与加权行报告按相应权重得到的效应估计；NA 表示该表项未单独列示。原始组均值用于计算朴素差异，调整后的估计由分层或加权公式计算。

朴素差异同时包含处理效果和客户选择的影响。按风险层标准化后，估计为 0.2500；按风险分数估计倾向得分并加权后，估计为 0.2199。在风险变量充分调整混杂且正值性成立时，这些结果可解释为优惠券提高留存率的效应估计。

倾向得分模型中，\texttt{risk\_score} 的系数为 1.0084，说明风险分数越高，越可能收到优惠券。平衡诊断显示，\texttt{risk\_score} 的原始 SMD 为 0.9710，加权后降为 0.2243；\texttt{support\_tickets\_90d} 的 SMD 从 1.0180 降为 0.3257。加权降低了观测协变量的不平衡程度，但残余差异仍较明显，管理结论需保留相应限制。

面板数据的 DID 结果为：
\begin{table}[htbp]
\centering
\caption{门店留存率的两期比较；各数值均为比例}
\begin{tabular}{lrrr}
\toprule
\rowcolor{BDMTableHead}
组别 & 干预前 & 干预后 & 变化\\
\midrule
处理门店 & 0.6100 & 0.7025 & 0.0925\\
对照门店 & 0.7125 & 0.7275 & 0.0150\\
\bottomrule
\end{tabular}
\end{table}
因此 DID 估计为
\[
  0.0925-0.0150=0.0775.
\]
这可以解释为：在用对照门店同期变化校正共同时间冲击后，处理门店的留存率额外提高约 7.75 个百分点。其识别依据是平行趋势、无预期影响和比较对象稳定等条件。

\section{因果研究报告}

因果估计进入管理决策前，应至少回答以下问题：
\begin{enumerate}
  \item 干预定义是否清晰，是否存在多个版本；
  \item 处理组和对照组在干预前是否可比；
  \item 哪些混杂因素已被观察，哪些重要因素可能遗漏；
  \item 是否存在足够重叠，是否有极端权重；
  \item DID 是否有干预前多期趋势支持；
  \item 是否存在溢出效应，例如客户之间转赠优惠券；
  \item 估计对象是 ATE、ATT 还是某个特定人群的局部效果；
  \item 统计估计是否足以支持行动成本和实施风险。
\end{enumerate}

因果研究报告应包含：业务问题、处理定义、结果定义、样本与时间窗口、识别策略、关键假设、估计结果、稳健性检查、不可识别风险、行动建议和后续实验计划。若结论用于将策略推广至目标人群，还应说明预期成本、容量约束、客户公平性和监控指标。

\section*{本章小结}
\addcontentsline{toc}{section}{本章小结}

预测模型可以发现风险对象，因果推断评估行动是否改变结果。潜在结果框架把处理效果写成反事实差异；随机实验提供清晰识别基准；观测性研究需要依赖一致性、SUTVA、条件可交换性和重叠性；倾向得分方法通过分层、匹配或加权改善可比性；DID 使用前后变化和对照组趋势校正共同冲击。所有因果估计都必须与识别假设、业务机制、数据质量和管理成本一起报告。

双重稳健估计的条件偏差由倾向得分误差和结果回归误差的乘积决定，交叉拟合分离
辅助函数估计与效应计算。分期处理还需明确首次采用时间、对照组和效应汇总权重。
评价企业平台的效果时，应区分整体上线与新增语义服务的目标效应，并把实施、培训和复核成本计入净收益。将效应估计转化为行动策略，还需明确收益、成本和资源限制，并据此建立优化模型，见第9章。

\chapterreading{潜在结果的基本思想见 Rubin 与 Holland\cite{rubin1974,holland1986}，
倾向得分的经典论述见 Rosenbaum 与 Rubin\cite{rosenbaum1983}。
Imbens 与 Rubin、Hernan 与 Robins 分别系统讨论设计、识别与估计
\cite{imbens2015,hernan2020}，因果图可参考 Pearl\cite{pearl2009}。
准实验的实证传统见 Angrist 与 Pischke\cite{angrist2009}。
双重机器学习和策略学习将灵活估计与目标参数联系起来\cite{dml08,policy08}；
这些方法的适用性取决于其识别、重叠和收敛条件。}

\section*{思考与习题}
\addcontentsline{toc}{section}{思考与习题}

\begin{exercise}[识别与稳健性]
逐步验证双重稳健分数在倾向模型正确时的条件期望。若关键设备健康变量未被记录，为什么两个预测模型都拟合良好仍不能保证因果识别？
\end{exercise}
\begin{exercise}[平台采用研究]
为一项分期部署的 AI 问数系统设计成效评估：写出处理时点、结果、候选对照、干扰与共同政策；说明如何避免把早期采用者的已有改善当成系统效果。
\end{exercise}

\begin{exercise}
围绕第8章样例数据完成下列任务：
\begin{enumerate}
  \item 计算优惠券处理组和对照组的朴素留存率差异；
  \item 按 \texttt{risk\_band} 分层，解释为什么分层估计与朴素估计不同；
  \item 运行倾向得分加权代码，并解释加权前后标准化均值差异的变化；
  \item 使用门店面板数据手工计算 DID；
  \item 写一段不超过 400 字的管理建议，说明是否应扩大优惠券策略，并列出至少三个识别风险。
\end{enumerate}
\end{exercise}

\begin{exercise}[标准化与目标总体]
重算表~\ref{tab:causal-example}~的组内和总体差异。如果目标总体中高难度订单占80\%，
分别计算两个行动的标准化延期率，说明平均效应为什么在本例中不变。再构造两个
难度层效应不同的例子，使目标总体构成改变时平均效应也改变。由此解释 ATE 与 ATT
为什么可能不同，而不把差异自动解释为某个估计方法错误。
\end{exercise}

\begin{exercise}[敏感性与行动强度]
一项干预的统计区间为 \([0.02,0.08]\)，结果越大越好，业务价值阈值为0.03。
分别在未测混杂偏差绝对值不超过0.01和0.04的假设下给出扩展范围。区分“方向可能
为正”“收益足以覆盖成本”和“值得扩大试点”三个判断，并提出一项可以减少识别
不确定性的研究设计。说明仅增加同类观测样本为何未必解决该问题。
\end{exercise}
